\uuid{48kB}
\exo7id{3556}
\titre{exo7 3556}
\auteur{quercia}
\organisation{exo7}
\datecreate{2010-03-10}
\isIndication{false}
\isCorrection{true}
\chapitre{Réduction d'endomorphisme, polynôme annulateur}
\sousChapitre{Polynôme annulateur}
\cpge{mpsi-mp2i.applications-lineaires.b, pcsi.applications-lineaires.b, ptsi.applications-lineaires.b, mp-mpi.reduction-endomorphismes.d, pc.reduction.c, psi.reduction-endomorphismes.c, pt.reduction-endomorphismes.c, mp-mpi.reduction-endomorphismes.g, pc.complements-espaces-vectoriels.d, psi.complements-espaces-vectoriels.d, mp-mpi.reduction-endomorphismes.i, pc.reduction.d, psi.reduction-endomorphismes.d}
\module{Algèbre}
\niveau{L2}
\difficulte{}

\contenu{
\question{
Soit $f$ un endomorphisme diagonalisable d'un ev~$E$ de dimension finie,
$\lambda$ une valeur propre de~$f$ et $p_\lambda$ le projecteur sur le
sous-espace propre associé parallèlement à la somme des autres sous-espaces
propres. Montrer que $p_\lambda$ est un polynôme en~$f$.
}
\reponse{
Soit $P$ un polynôme tel que $P(\lambda) = 1$ et $P(\mu) = 0$ pour
toutes les autres valeurs propres, $\mu$, de~$f$. Alors $p_\lambda = P(f)$.
}
}
